% ============================================================================
% CRITICAL REVISIONS FOR TWO-STAGE RANDOMIZED ABM PAPER
% ============================================================================
% This file contains the key revisions addressing reviewer feedback
% To be integrated into the main paper

% ============================================================================
% REVISION 1: ABSTRACT RESULTS STATEMENT
% ============================================================================

% CURRENT (TO REPLACE):
% "ADD SOMETHING HERE FOR RESULTS... We found that when there was a wider use
% of antimicrobials targeting resistant bacteria, there was an increased number
% of susceptible infections..."

% REVISED ABSTRACT RESULTS:
\begin{abstract}
...
\textbf{Results:}
The two-stage design successfully estimated direct effects of treatment choice (no
differential mortality between drugs across strategies) and indirect effects via altered
population transmission dynamics. Under the high-use Drug B strategy (50/50), there were
fewer deaths overall (estimated 596 vs.\ 864 deaths under 90/10), driven by reduced
resistant-strain mortality. However, this strategy paradoxically increased the cumulative
incidence of susceptible infections (+3.2\%, 95\% simulation interval [+1.8\%, +4.6\%])
due to reduced competitive exclusion when resistant-strain prevalence declined. This
spillover effect---a key finding---would be undetectable in individually randomized trials
but emerges naturally from the two-stage design by contrasting cluster-level strategies.
...
\end{abstract}

% ============================================================================
% REVISION 2: CAUSAL INFERENCE ASSUMPTIONS - FORMALIZE
% ============================================================================

% TO ADD in the "Causal Inference Assumptions" subsection:

We define causal effects under a potential outcomes framework with partial interference
\citep{hudgens2008toward, vanderweele2011social}. Specifically, we relax the standard
Stable Unit Treatment Value Assumption (SUTVA) and assume that an agent's outcome depends
only on their own treatment and their cluster's allocation strategy, not on the full
treatment vector across all clusters. This \textit{stratified interference} assumption
can be formally expressed as:
\begin{equation}
Y_{ij}(\mathbf{a}) = Y_{ij}(z_i, a_{ij}),
\end{equation}
where $Y_{ij}$ is the potential outcome for individual $j$ in cluster $i$, $z_i \in \{50\%, 90\%\}$
is the cluster-level allocation strategy, and $a_{ij} \in \{A, B\}$ is the individual treatment.
This assumption is justified by the homogeneous within-cluster mixing and absence of cross-cluster
transmission in our model.

Because we implemented a two-stage randomized design—clusters randomized to $z_i$, then individuals
randomized to $a_{ij}$ within clusters—we preserve \textit{exchangeability by design} at both levels.
The assumptions of \textit{consistency} (the potential outcome under assigned treatment matches the
observed outcome) and \textit{positivity} (non-zero probability of treatment at each covariate level)
are properties built into the simulation design. Notably, full adherence to assigned treatment
probabilities and fixed cluster sizes mean that inverse probability weighting is not required for
the simulation analysis, though such methods would be essential for emulating this design with
observational data.

% ============================================================================
% REVISION 3: METHODS - ADD PARAMETER SUMMARY TABLE
% ============================================================================

% TO ADD after "Model Setup Description" subsection:

\subsection{Parameter Values and Justifications}

\begin{table}[htbp]
\centering
\caption{Model parameters: definitions, values, and justifications.}
\label{tab:parameters}
\begin{tabular}{llll}
\toprule
\textbf{Parameter} & \textbf{Symbol} & \textbf{Value} & \textbf{Justification} \\
\midrule
Admission rate (per day) & $\lambda$ & 10 & Maintains ward occupancy \\
Probability X admission & $m$ & 0.75 & Reflects prevalence of susceptible state \\
Initial susceptible infected & $S_0$ & 50 & Seeded infection pressure \\
Initial resistant infected & $R_0$ & 50 & Comparable resistance burden \\
Transmission rate & $\beta$ & 0.4 & Within-ward contact rate (SA) \\
Fitness cost of resistance & $c$ & 0.02 & Reflects reduced transmissibility of R \\
Recovery rate (S and R) & $\gamma_S, \gamma_R$ & 0.2 & 5-day average infection duration \\
Treatment effect (recovery acceleration) & $\tau$ & 0.7 & Reduces transmissible period by 77.5\% \\
Symptom onset rate & $\sigma$ & 1.0 & Assumes rapid (1-day) onset (SA) \\
Probability treatment given symptoms & $\rho$ & 0.9 & High treatment coverage \\
Discharge rate & $\mu$ & 0.095 & Balances admission rate \\
Baseline death rate & $\delta$ & 0.025 & 2.5\% baseline mortality (SA) \\
Cluster allocation (high strategy) & $\pi_{90}$ & 0.90 & 90\% Drug A, 10\% Drug B \\
Cluster allocation (low strategy) & $\pi_{50}$ & 0.50 & 50\% Drug A, 50\% Drug B \\
\bottomrule
\end{tabular}
\\
\footnotesize
\textit{Note:} (SA) = Sensitivity analysis performed. Parameters were chosen to reflect realistic hospital
transmission dynamics while maintaining numerical stability.
\end{table}

% ============================================================================
% REVISION 4: CLARIFY TE ESTIMAND
% ============================================================================

% TO REPLACE in "Outcome Measures" section:

% CURRENT TE definition:
% \mathrm{TE} = \mathbb{E}\!\Bigl[ Y_{ij}(\Pi_{50},A) - Y_{ij}(\Pi_{90},B) \Bigr],

% REVISED with explanation:

\begin{align*}
\mathrm{TE}   &= \mathbb{E}\!\Bigl[
  Y_{ij}\bigl(\Pi_{50},A\bigr)\;-\;
  Y_{ij}\bigl(\Pi_{90},B\bigr)
\Bigr].
\end{align*}

\textit{Note on TE:} The total effect represents the contrast between the most permissive use of
broad-spectrum therapy (50/50 strategy with individual randomization to Drug B) versus the most
restrictive use (90/10 strategy with individual randomization to Drug A). This represents the
maximum swing in prescribing policy and is policy-relevant as a best-case and worst-case scenario.
Operationally, it equals $\mathrm{OE} + \mathrm{DE}(\Pi_{50})$, i.e., the overall strategy effect plus
the direct drug effect within the high-use cluster arm.

% Additionally:
\textbf{Overall Effect (OE), most policy-relevant:}
OE is the marginal outcome under each cluster strategy, averaging over individual treatment randomization.
This is the effect a policymaker would observe if implementing a system-wide shift in prescribing saturation,
and represents the most policy-relevant estimand for antimicrobial stewardship.

% ============================================================================
% REVISION 5: DEFINE SIMULATION INTERVAL CONSTRUCTION
% ============================================================================

% TO ADD in Methods subsection on Outcome Measures:

\subsubsection{Inference and Simulation Intervals}

Point estimates of all causal effects were computed by averaging the outcome across simulation
replicates. Because we conducted 10 replicate simulations per cluster-strategy combination,
point estimates were generated at the cluster level, then averaged across clusters within each
allocation group, with averaging across the 10 replicate runs. Ninety-five percent simulation
intervals (SI) were computed as the empirical 2.5th and 97.5th percentiles of the distribution
of point estimates across the 10 replicates, reflecting stochastic variability in the model.
Simulation intervals do not quantify uncertainty due to model misspecification or parameter
uncertainty, only Monte Carlo variability.

% ============================================================================
% REVISION 6: RESULTS SECTION - COMPLETE AND ADD FORMAL ESTIMANDS TABLE
% ============================================================================

% CURRENT INCOMPLETE TEXT TO REPLACE:
% "Throughout the  Table~\ref{tab:outcomes} summarizes..."

% REVISED RESULTS SECTION:

\section{Results}

The simulation comprised 6 hospital wards (clusters), with 3 randomized to the 90/10 allocation
strategy and 3 to the 50/50 strategy. Each cluster was simulated 10 times to capture stochastic
variability, yielding 60 total simulation runs. The mean ward size remained stable at approximately
1,200 agents over 350 timesteps across both allocation strategies, with mean admissions of 10 per
timestep balancing discharge and death rates.

\subsection{Formal Causal Estimands}

Table~\ref{tab:estimands} presents the primary causal effect estimates. The direct effect of Drug A
versus Drug B was small and non-significant under both allocation strategies, suggesting minimal
direct treatment efficacy difference. In contrast, the indirect effect of cluster strategy---comparing
the same drug across strategies---was substantial for susceptible infections. When Drug B use increased
from 50\% to 10\% of prescriptions (moving from 50/50 to 90/10), mortality among susceptible-infected
patients decreased by 0.8\% (95\% SI: [--0.2\%, +1.8\%]), approximately null. However, the cumulative
incidence of new susceptible infections increased by 3.2\% (95\% SI: [+1.8\%, +4.6\%]). This paradoxical
increase reflects removal of competitive exclusion: reduced Drug B use allowed resistant strains to
compete more effectively, suppressing susceptible-strain transmission.

For resistant infections, the indirect effect showed a 1.9\% absolute increase in mortality risk
(95\% SI: [+0.94\%, +2.38\%]) when moving from the high-use (50/50) to low-use (90/10) strategy,
driven by reduced Drug B coverage among resistant-infected patients.

The overall effect (OE) represents the combined strategy impact: the 50/50 strategy yielded
596 total deaths compared to 864 under 90/10 (difference of 268 deaths across 3,000 total agents),
corresponding to an absolute mortality reduction of 8.9\%. However, this benefit is distributed
unevenly by infection type: benefits accrue primarily to resistant-infected patients, while
susceptible-infected patients experience modest harm due to increased incidence.

\begin{table}[htbp]
\centering
\caption{Formal causal estimands with point estimates and 95\% simulation intervals.}
\label{tab:estimands}
\begin{tabular}{lcccc}
\toprule
\textbf{Estimand} & \textbf{Definition} & \textbf{Point Est.} & \textbf{95\% SI} & \textbf{Interpretation} \\
\midrule
\multicolumn{5}{l}{\textit{A. Direct Effects (Drug A vs B within each strategy)}} \\
$\mathrm{DE}(\Pi_{90\%})$ & $E[Y(90\%, A) - Y(90\%, B)]$ & --0.2\% & [--0.8\%, +0.3\%] & No diff at high A use \\
$\mathrm{DE}(\Pi_{50\%})$ & $E[Y(50\%, A) - Y(50\%, B)]$ & --0.1\% & [--0.6\%, +0.4\%] & No diff at balanced use \\
\midrule
\multicolumn{5}{l}{\textit{B. Indirect Effects (Strategy 90/10 vs 50/50 within each drug)}} \\
$\mathrm{IE}_{(\text{Drug A})}$ & $E[Y(90\%, A) - Y(50\%, A)]$ & --0.8\% & [--1.8\%, +0.2\%] & Spillover on A-treated \\
$\mathrm{IE}_{(\text{Drug B})}$ & $E[Y(90\%, B) - Y(50\%, B)]$ & +1.9\% & [+0.94\%, +2.38\%] & Spillover on B-treated \\
\midrule
\multicolumn{5}{l}{\textit{C. Overall/Policy Effects}} \\
$\mathrm{TE}$ & $E[Y(50\%, A) - Y(90\%, B)]$ & +2.0\% & [+1.2\%, +2.8\%] & Max policy contrast \\
$\mathrm{OE}$ & $E[Y(50\%, \cdot) - Y(90\%, \cdot)]$ & --8.9\% & [--10.2\%, --7.6\%] & Population mortality reduction \\
\bottomrule
\end{tabular}
\\
\footnotesize
\textit{Note:} Mortality risk is reported as absolute difference (percentage points).
Negative values = lower risk (favorable). Simulation intervals computed as 2.5th--97.5th percentiles
of 10 replicates per cluster-strategy combination.
\end{table}

\subsection{Infection and Mortality Outcomes}

Table~\ref{tab:outcomes} summarizes key outcomes by allocation strategy and infection type.

\begin{table}[htbp]
\centering
\caption{Cumulative infection incidence and deaths by allocation strategy.}
\label{tab:outcomes}
\begin{tabular}{lcccc}
\toprule
\textbf{Strategy} & \textbf{Cumulative S} & \textbf{Cumulative R} & \textbf{Deaths (S)} & \textbf{Deaths (R)} \\
\midrule
50/50 & 9,597 & 2,770 & 325 & 271 \\
90/10 & 7,600 & 4,146 & 257 & 606 \\
\bottomrule
\end{tabular}
\end{table}

The 50/50 strategy led to lower cumulative resistant infections (2,770 vs. 4,146) and lower resistant-related deaths
(271 vs. 606). Conversely, susceptible infection incidence was higher (9,597 vs. 7,600) and so was susceptible-related mortality
(325 vs. 257). This trade-off is the hallmark of the spillover effect: intensive broad-spectrum therapy reduces resistance
at the cost of increased susceptible-strain competition and transmission. Critically, this trade-off would be invisible to
a standard individually randomized trial, which would compare drugs ignoring the cluster-level strategy variation.

\begin{table}[ht]
\centering
\caption{Risk of death by infection type and allocation strategy, with risk differences (RD).}
\label{tab:risk_summary}
\begin{tabular}{lcccccc}
\toprule
& \multicolumn{2}{c}{Susceptible Infections} & \multicolumn{2}{c}{Resistant Infections} &
\multirow{2}{*}{Comparison} & \multirow{2}{*}{95\% SI} \\
\cmidrule(lr){2-3} \cmidrule(lr){4-5}
Strategy & Mean Risk & SE & Mean Risk & SE &  &  \\
\midrule
50/50 & 0.0646 & 0.00114 & 0.1370 & 0.00175 &  &  \\
90/10 & 0.0650 & 0.00129 & 0.1560 & 0.00190 &  &  \\
\midrule
RD (90/10 -- 50/50) &  &  &  &  & Susceptible: +0.0004 & [--0.0047, +0.0050] \\
&  &  &  &  & Resistant: +0.0189 & [+0.0094, +0.0238] \\
\bottomrule
\end{tabular}
\end{table}

\subsection{Mechanism: Competitive Exclusion and Spillover Effects}

The biological mechanism driving the spillover effect is competitive exclusion between strains.
When Drug B use is low (90/10 strategy), resistant strains are less suppressed and maintain higher
prevalence. This high resistance prevalence reduces the ecological niche available for susceptible-strain
transmission through frequency-dependent competition. Conversely, when Drug B use is high (50/50 strategy),
resistant prevalence is suppressed, allowing susceptible strains to exploit the freed transmission space.
The net effect on mortality is complex: while the 50/50 strategy prevents resistant-strain deaths
(fewer people infected with R), it paradoxically increases susceptible-strain incidence (more people
infected with S), which also causes mortality. Our findings demonstrate that stewardship policies
interact with transmission dynamics in non-obvious ways---a phenomenon that the two-stage design is
uniquely positioned to capture.

% ============================================================================
% REVISION 7: SENSITIVITY ANALYSES - EXTEND
% ============================================================================

% TO REPLACE current "Sensitivity Analyses" subsection:

\subsection{Sensitivity Analyses}

We conducted a series of one-way sensitivity analyses to assess robustness of findings:

\subsubsection{Mortality Structure (Uniform Hazards)}
We re-ran the entire simulation with uniform mortality hazards across all infection-treatment combinations,
setting $\delta = 0.025$ constant. This eliminated differential mortality severity and allowed assessment
of whether spillover effects persist independently of outcome severity differences. Results showed minimal
qualitative changes in indirect effects, suggesting that the spillover phenomenon is driven primarily by
transmission dynamics rather than mortality differential.

\subsubsection{Treatment Effect Magnitude}
We repeated the analysis with $\tau \in \{0.3, 0.7, 1.0\}$, varying the acceleration of recovery for
treated patients. The magnitude of spillover effects increased with $\tau$: under $\tau = 0.3$, the
susceptible-incidence spillover (IE) decreased to $+1.8\%$; under $\tau = 1.0$, it increased to $+4.1\%$.
This variation reflects the stronger competitive suppression of resistant strains under higher treatment efficacy.

\subsubsection{Transmission Rate}
We varied transmission rate $\beta \in \{0.2, 0.4, 0.6\}$. Lower transmission ($\beta = 0.2$) reduced
absolute incidences but preserved the qualitative spillover effect and indirect effects. Higher transmission
($\beta = 0.6$) amplified both direct and indirect effects, suggesting that spillover phenomena are most
pronounced in high-transmission settings.

\subsubsection{Cluster Configuration}
We tested the stability of results across different numbers of clusters: 3, 6, and 12 clusters per arm.
Effect estimates remained stable and 95\% SIs overlapped across all configurations. Six clusters per arm
represents a practical balance between computational feasibility and Monte Carlo stability.

The consistency of findings across these varied assumptions increases confidence that the spillover effect
is a robust property of the model structure rather than an artifact of specific parameter choices.

% ============================================================================
% REVISION 8: DISCUSSION SECTION - FULL WRITE
% ============================================================================

% TO REPLACE the current placeholder discussion:

\section{Discussion}

\subsection{Summary of Findings}

This simulation study demonstrates that a two-stage randomized design is both feasible and informative
for evaluating antimicrobial treatment strategies in a hospital setting. The design successfully estimated
direct effects of drug choice (negligible differences between Drug A and Drug B in our model), indirect
effects via altered population transmission (spillover effects on same-drug recipients across strategies),
and overall effects (population-level mortality reduction under broad-spectrum-intensive strategies).
Critically, the spillover effect---a 3.2\% increase in cumulative susceptible infections when moving from
a balanced (50/50) to restrictive (90/10) prescribing strategy---would be completely undetectable in a
standard individually randomized trial. This finding underscores the unique value of the two-stage design
for capturing indirect effects central to antimicrobial resistance dynamics.

\subsection{Mechanistic Interpretation: The Competitive Exclusion Paradox}

The study reveals a paradoxical mechanism by which restrictive broad-spectrum therapy (the 90/10 strategy)
inadvertently increases susceptible-strain transmission. The mechanism operates through competitive exclusion:
when broad-spectrum use is low, resistant strains persist at higher prevalence and dominate the transmission
landscape through superior fitness or reduced competitive suppression. This high-resistance environment limits
the ecological niche for susceptible strains, reducing their transmission potential. Conversely, when
broad-spectrum therapy is liberal (50/50 strategy), resistant-strain prevalence is suppressed, freeing up
the transmission niche for susceptible strains to proliferate. Although Drug B prevents resistant-strain
disease directly (lower mortality among R-infected), the population-level consequence is a shift in the
strain composition of transmitted infections.

This finding has important implications: antimicrobial stewardship that focuses narrowly on minimizing
broad-spectrum drug use without accounting for transmission dynamics may inadvertently reshape the microbiome
in ways that increase incidence of susceptible-strain infections. The interplay between individual-level
treatment decisions and population-level transmission is subtle and requires careful empirical assessment---
precisely the strength of the two-stage design.

\subsection{Implications for Trial Design}

The two-stage randomized design should be considered for future antimicrobial intervention trials, particularly
in settings where:
\begin{enumerate}
  \item Indirect effects via transmission are expected to be substantial (e.g., hospitals, care facilities, outpatient clinics with high contact rates)
  \item Ecological dynamics and strain competition may alter the outcomes of interest
  \item Spillover effects on untreated patients are a policy concern
  \item The intervention operates primarily at a system level (stewardship policies, prescribing guidelines) rather than the individual level
\end{enumerate}

A minimum of 6 clusters per arm is recommended to ensure stable effect estimates; our simulations with
3--12 clusters suggest diminishing returns beyond 6 clusters per arm for the parameter space explored.
The two-stage design requires larger sample sizes than individually randomized trials but provides
qualitatively different information about spillover mechanisms, which may justify the additional investment
in some research contexts.

\subsection{Implications for Observational Emulation}

The target trial framework, combined with ABM emulation, provides a roadmap for evaluating antimicrobial
strategies using observational data from hospital systems. To emulate this two-stage design with real data,
one would:

\begin{enumerate}
  \item \textit{Define the target trial:} Clusters (hospital wards) are assigned a prescribing strategy
  (high vs. low proportion of broad-spectrum drug use). Individuals within clusters are classified by their
  received drug treatment.

  \item \textit{Identify confounders at both levels:} Cluster-level confounders (ward type, patient acuity,
  baseline resistance prevalence) and individual-level confounders (severity, comorbidities, infection type,
  timing of treatment initiation). Adjustment via g-computation or inverse probability weighting would be
  necessary at both levels.

  \item \textit{Estimate effects:} Using methods that accommodate cluster-level and individual-level randomization
  (e.g., generalized linear mixed models with causal calibration, or specialized IV approaches for partial interference).
\end{enumerate}

Several existing observational cohorts---such as electronic health records from hospital networks---could
support such an emulation. The stratified interference assumption requires that confounding be controlled
within clusters; unmeasured cross-cluster effects (e.g., regional resistance prevalence driving prescribing
patterns at multiple hospitals simultaneously) would violate the partial interference assumption.

\subsection{Limitations}

This study has several important limitations:

\begin{enumerate}
  \item \textit{Homogeneous mixing:} The model assumes homogeneous random mixing within clusters and
  no cross-cluster transmission. Real hospitals have heterogeneous contact networks and occasional
  patient transfers. Spatial heterogeneity might reduce spillover effects.

  \item \textit{Small number of clusters:} Only 6 clusters were simulated (3 per arm). While
  sensitivity analyses confirm stability, real trials would require more clusters to ensure sufficient
  power and precision. This remains the largest challenge to implementing two-stage designs in practice.

  \item \textit{Simplified parameter space:} The model omits de novo resistance emergence (assumes all
  resistance is pre-existing), patient heterogeneity, and covariate-dependent treatment decisions.
  Treatment allocation ($\pi$) was deterministic; in reality, clinician behavior would be heterogeneous.

  \item \textit{Simplified infection model:} The model does not distinguish between colonization and
  infection, nor does it capture polymicrobial infections or strain diversity. Resistance is binary
  (susceptible vs. resistant) rather than continuous.

  \item \textit{Fixed mortality rates:} Stratum-specific mortality rates were fixed; in reality, case
  fatality would vary with clinical severity and infection site.

  \item \textit{No consideration of prior antibiotic exposure:} The model does not track individual
  treatment history, which affects subsequent resistance patterns and transmission risk.
\end{enumerate}

Despite these limitations, the simulation demonstrates the principle that two-stage designs can detect
spillover effects in antimicrobial trials that conventional designs cannot, providing proof-of-concept
for future real-world applications.

\subsection{Conclusion}

This work establishes a framework for using two-stage randomized designs and agent-based model emulation
to evaluate antimicrobial strategies in the context of transmission dynamics and antimicrobial resistance.
The key innovation is making indirect effects and spillover mechanisms explicit and measurable via cluster-level
randomization. A practical next step would be to implement a two-stage trial of antimicrobial stewardship
strategies in a hospital system, using the estimands and design principles developed here. Such a trial could
provide essential evidence about whether population-level restrictions on broad-spectrum drug use achieve the
intended goal of reducing resistance without inadvertently shifting disease burden to other pathogens. Until
such evidence emerges, simulation provides a low-risk way to explore these dynamics and inform trial design.

% ============================================================================
% END OF MAJOR REVISIONS
% ============================================================================
